\documentclass[12pt,a4paper]{article}
\usepackage[margin=25mm, headheight=15pt, headsep=10pt]{geometry}
\usepackage{fontspec}
\usepackage[table]{xcolor} % Note: table option is required for rowcolors
\usepackage{titlesec}
\usepackage{tocloft}
\usepackage{fancyhdr}
\usepackage{booktabs}
\usepackage{tcolorbox}
\tcbuselibrary{skins, breakable}
\usepackage[colorlinks=true, linkcolor=emerald, urlcolor=emerald, bookmarks=true]{hyperref}

\definecolor{emerald}{RGB}{0,102,51}
\definecolor{darkred}{RGB}{139,0,0}
\definecolor{lightgray}{RGB}{245,245,245}

\usepackage[nil]{babel}
\babelprovide[import, main]{bengali}
\babelprovide[import]{arabic}
\babelfont{rm}[Renderer=HarfBuzz,Scale=1.0]{Noto Serif Bengali}
\babelfont{sf}[Renderer=HarfBuzz]{Noto Sans Bengali}
\babelfont{tt}[Renderer=HarfBuzz]{Noto Sans Bengali}
\babelfont[arabic]{rm}[Renderer=HarfBuzz,Scale=1.1]{Noto Naskh Arabic}

% Section styling
\titleformat{\section}{\Large\bfseries\color{emerald}}{\thesection.}{0.5em}{}
\titleformat{\subsection}{\large\bfseries\color{emerald!85!black}}{\thesubsection.}{0.5em}{}
\titleformat{\subsubsection}{\normalsize\bfseries\color{emerald!70!black}}{\thesubsubsection.}{0.5em}{}
\titlespacing*{\section}{0pt}{18pt}{8pt}

% Fancy Header/Footer setup
\pagestyle{fancy}
\fancyhf{} % clear all header and footer fields
\fancyhead[L]{\color{emerald}\small\textbf{\nouppercase{\leftmark}}}
\fancyfoot[L]{\scriptsize\color{black!50}{\lat{AI}}-সংকলিত, আলেম-যাচাইকৃত নয়}
\fancyfoot[C]{\color{emerald!70!black}\thepage}
\renewcommand{\headrulewidth}{0.4pt}
\renewcommand{\headrule}{\hbox to\headwidth{\color{emerald}\leaders\hrule height \headrulewidth\hfill}}
\renewcommand{\footrulewidth}{0pt}

% Custom boxes
% 1. Ayat/Hadith Box
\newtcolorbox{ayatbox}[1][]{
    enhanced, breakable, 
    colback=emerald!5, 
    colframe=emerald, 
    boxrule=0.5pt, 
    arc=3pt, 
    left=10pt, right=10pt, top=10pt, bottom=10pt,
    #1
}

% 2. Quote/Translation Box
\newtcolorbox{quotebox}[1][]{
    enhanced, breakable,
    frame hidden,
    colback=lightgray,
    borderline west={3pt}{0pt}{emerald},
    left=15pt, right=10pt, top=10pt, bottom=10pt,
    sharp corners,
    #1
}

% 3. AI-generated notice box
\newtcolorbox{warnbox}[1][]{
    enhanced, breakable,
    colback=red!4,
    colframe=darkred,
    boxrule=0.8pt,
    arc=3pt,
    left=10pt, right=10pt, top=10pt, bottom=10pt,
    #1
}

% Commands
\newcommand{\arab}[1]{\foreignlanguage{arabic}{#1}}
\newcommand{\kib}[1]{\textcolor{emerald}{\textbf{#1}}}

% Latin (English) fallback: Noto Serif Bengali has NO Latin glyphs
\newfontfamily\latfont[Renderer=HarfBuzz]{Noto Serif}
\newcommand{\lat}[1]{{\latfont #1}}

\setlength{\parskip}{5pt}
\setlength{\parindent}{0pt}

% Fix for missing bullet in Noto Serif Bengali
\renewcommand{\labelitemi}{$\bullet$}



\begin{document}
\begin{center}{\Large\bfseries\color{emerald} হাদিস ও আসার — ছবি আঁকার বিধান (তাসবীর)}\end{center}
\vspace{1em}
\section*{হাদিস ও আসার — ছবি আঁকার বিধান (তাসবীর)}

\begin{warnbox}
 \textbf{গুরুত্বপূর্ণ নোটিশ (\lat{Important} \lat{Notice}):} এই কন্টেন্টটি \lat{AI} (কৃত্রিম বুদ্ধিমত্তা) দিয়ে প্রস্তুত ও সংকলিত। \lat{AI} কঠোর সূত্র-মান নীতি মেনে শুধু নির্ভরযোগ্য উৎস থেকে তথ্য সংগ্রহ করেছে — কেবল সহীহ/হাসান হাদিস ও সহীহ সনদের আসার রাখা হয়েছে, দুর্বল সনদ বাদ দেওয়া হয়েছে। তবে এটি কোনো আলেম বা মুহাদ্দিস কর্তৃক যাচাইকৃত নয়।
\end{warnbox}

\begin{quote}
\textbf{পড়ার নিয়ম:} প্রতিটি হাদিসের সাথে থাকবে — \textbf{সূত্র কার্ড} (সনদের মান) $\rightarrow$ \textbf{পটভূমি} $\rightarrow$ \textbf{পূর্ণ বর্ণনা (বাংলা, \lat{pure})} $\rightarrow$ \textbf{ধাপে ধাপে বিশ্লেষণ} $\rightarrow$ \textbf{শব্দের ব্যাখ্যা} $\rightarrow$ \textbf{গুরুত্বপূর্ণ শিক্ষা}। \\
\textbf{মূল নীতিমালা:} শুধুমাত্র সহীহ/হাসান হাদিস এবং সহীহ সনদের আসার (আথার) গ্রহণ করা হয়েছে। দুর্বল সনদ সম্পূর্ণ বাদ দেওয়া হয়েছে।
\end{quote}

\medskip\hrule\medskip

\subsection*{পার্ট ১ — ছবি নির্মাণের নিষেধ ও সতর্কবার্তা}

\subsubsection*{হাদিস ১ — সবচেয়ে কঠিন শাস্তি ছবি নির্মাতাদের (সহীহ বুখারি ৫৯৫০, সহীহ মুসলিম ২১০৯)}

\#\#\#\# ১. সূত্র কার্ড

\begin{table}[h]\centering\small
\begin{tabular}{ll}
বিষয় & বিবরণ \\
\hline
\textbf{বর্ণনাকারী} & আবদুল্লাহ ইবনে মাসউদ (রাঃ) \\
\textbf{গ্রন্থ} & সহীহ বুখারি ৫৯৫০; সহীহ মুসলিম ২১০৯; মুসনাদ আহমাদ ৩৫৫৮ \\
\textbf{অধ্যায়} & কিতাবুল লিবাস — বাবু আজাবিল মুসাওয়িরীন ইয়াওমাল কিয়ামাহ (ছবি নির্মাতাদের শাস্তির অধ্যায়) \\
\textbf{সনদের মান} & \textbf{সহীহ} (বুখারি-মুসলিম) \\
\hline
\end{tabular}
\end{table}

\#\#\#\# ২. পটভূমি (\lat{Context})

ছবি নিয়ে ইসলামের সবচেয়ে বিখ্যাত সতর্কবার্তা। মাসরূক (রাঃ) এক বাড়ির আঙিনায় ছবি দেখে এই হাদিসটি স্মরণ করেন। খেয়াল করুন — হাদিসে বলা হয়েছে \lat{“}আল-মুসাওয়িরুন\lat{”} (ছবি নির্মাতারা), শুধু \lat{“}মানুষের ছবি নির্মাতারা\lat{”} নয়। শব্দটি \textbf{সাধারণ (\lat{general})} — যে কেউ প্রাণবান সৃষ্টির ছবি বানায়।

\#\#\#\# ৩. পূর্ণ বর্ণনা (বাংলা, \lat{pure})

\begin{quote}
রাসূলুল্লাহ (সাল্লাল্লাহু আলাইহি ওয়া সাল্লাম) বলেছেন: \\
\textbf{\lat{“}নিশ্চয়ই কিয়ামতের দিন আল্লাহর কাছে মানুষের মধ্যে সবচেয়ে কঠিন শাস্তি হবে ছবি নির্মাতাদের।\lat{”}}
\end{quote}

\#\#\#\# ৪. ধাপে ধাপে বিশ্লেষণ

\textbf{ধাপ ১ — শব্দচয়ন:} \lat{“}আল-মুসাওয়িরুন\lat{”} (\arab{المصورون}) — ছবি বানানোর সঙ্গে জড়িত সবাই; মানুষের, পশু-পাখির, যেকোনো প্রাণবান সৃষ্টির।

\textbf{ধাপ ২ — শাস্তির গাম্ভীর্য:} \lat{“}সবচেয়ে কঠিন শাস্তি\lat{”} — এই স্তরের সতর্কবাণী খুব কম বিষয়ে এসেছে। কারণ — ছবি বানানো আল্লাহর \textbf{সৃষ্টির অনুকরণে (\lat{imitation} \lat{of} \lat{creation})} প্রবেশ।

\textbf{ধাপ ৩ — ব্যাপ্তি:} কেউ যদি মনে করে \lat{“}মানুষের ছবি নিষিদ্ধ, পশুপাখি জায়েজ\lat{”} — এই হাদিসের সাধারণ শব্দ তার বিপরীত। ইবনে আব্বাস, ইমাম নববি, খাত্তাবিসহ প্রায় সবাই এতে পশু-পাখি-মানুষ সব \textbf{প্রাণবান (\lat{animate}) সৃষ্টি} অন্তর্ভুক্ত বলেছেন।

\#\#\#\# ৫. শব্দের ব্যাখ্যা (\lat{Vocabulary})

\begin{table}[h]\centering\small
\begin{tabular}{ll}
শব্দ & অর্থ \\
\hline
\textbf{আশাদ্দুন নাসি (\arab{أشد} \arab{الناس})} & মানুষের মধ্যে সবচেয়ে কঠোর/কঠিন \\
\textbf{আযাবান (\arab{عذابا})} & শাস্তি \\
\textbf{আল-মুসাওয়িরুন (\arab{المصورون})} & ছবি নির্মাতারা — যারা প্রাণবান সৃষ্টির ছবি/মূর্তি বানায় \\
\hline
\end{tabular}
\end{table}

\#\#\#\# ৬. গুরুত্বপূর্ণ শিক্ষা (\lat{Key} \lat{Lessons})

\begin{enumerate}
\item \textbf{নিষেধ মানুষে সীমাবদ্ধ নয়:} \lat{“}ছবি নির্মাতা\lat{”} শব্দটি মানুষের ছবিকে একচেটিয়া করে না — পশু-পাখিসহ যেকোনো প্রাণবান সৃষ্টির ছবি এর অন্তর্ভুক্ত।
\item \textbf{সতর্কবার্তার গুরুত্ব:} ছবির বিষয়টি হালকা নয়; এটি ইসলামের সবচেয়ে গুরুত্বপূর্ণ সতর্কবার্তাগুলোর একটি।
\item \textbf{ইন্তিসাবের প্রশ্ন:} যিনি ছবি বানান (আঁকেন/গড়েন), তিনিই দায়ী — ছবিটা ঘরে রাখা বা না রাখা আলাদা মাসআলা।
\end{enumerate}
\medskip\hrule\medskip

\subsubsection*{হাদিস ২ — \lat{“}যা সৃষ্টি করেছ তাতে প্রাণ দাও\lat{”} (সহীহ বুখারি ৫৯৫১, সহীহ মুসলিম ২১০৭)}

\#\#\#\# ১. সূত্র কার্ড

\begin{table}[h]\centering\small
\begin{tabular}{ll}
বিষয় & বিবরণ \\
\hline
\textbf{বর্ণনাকারী} & আবদুল্লাহ ইবনে উমার (রাঃ) \\
\textbf{গ্রন্থ} & সহীহ বুখারি ৫৯৫১; সহীহ মুসলিম ২১০৭ \\
\textbf{অধ্যায়} & কিতাবুল লিবাস — বাবু আজাবিল মুসাওয়িরীন \\
\textbf{সনদের মান} & \textbf{সহীহ} (বুখারি-মুসলিম) \\
\hline
\end{tabular}
\end{table}

\#\#\#\# ২. পটভূমি (\lat{Context})

এটি ছবি নিষেধের সবচেয়ে শক্তিশালী বক্তব্য — কেন ছবি বানানো নিষিদ্ধ তার \textbf{কারণ (\lat{reason})} এখানেই: ছবি বানানো এমন এক দাবি যা কেউ পূরণ করতে পারে না — প্রাণ দেওয়া। ইবনে উমার (রাঃ) নবীজির থেকে সরাসরি শুনেছেন।

\#\#\#\# ৩. পূর্ণ বর্ণনা (বাংলা, \lat{pure})

\begin{quote}
\textbf{\lat{“}যারা এই ছবিগুলো বানায়, কিয়ামতের দিন তাদের শাস্তি দেওয়া হবে এবং তাদের বলা হবে — তোমরা যা সৃষ্টি করেছ, তাতে প্রাণ দাও।\lat{”}}
\end{quote}

\#\#\#\# ৪. ধাপে ধাপে বিশ্লেষণ

\textbf{ধাপ ১ — \lat{“}সৃষ্টি\lat{”} শব্দের প্রভাব:} এখানে ছবি বানানোকে \lat{“}সৃষ্টি\lat{”} বলা হয়েছে — অর্থাৎ এটি আল্লাহর একক বৈশিষ্ট্যের \textbf{চ্যালেঞ্জ (\lat{challenge})}।

\textbf{ধাপ ২ — প্রাণ দেওয়ার অসম্ভবতা:} \lat{“}তাতে প্রাণ দাও\lat{”} — মানুষ কখনো পারবে না। শাস্তির মর্মার্থ — ছবি নির্মাতা এমন দাবি করেছে যা সে পূরণে অক্ষম।

\textbf{ধাপ ৩ — নিষেধের ব্যাপ্তি:} আবারও লক্ষ্য করুন — \lat{“}এই ছবিগুলো\lat{”} (\arab{هذه} \arab{الصور})। হাদিসে মানুষ-পশু আলাদা করা হয়নি; প্রাণবান সৃষ্টির ছবিই উদ্দেশ্য।

\#\#\#\# ৫. শব্দের ব্যাখ্যা (\lat{Vocabulary})

\begin{table}[h]\centering\small
\begin{tabular}{ll}
শব্দ & অর্থ \\
\hline
\textbf{সুওয়ার (\arab{صور})} & ছবি — এখানে প্রাণবান সৃষ্টির ছবি \\
\textbf{ইয়াসনাউন (\arab{يصنعون})} & বানানো — তৈরি করা \\
\textbf{আহয়ু (\arab{أحيوا})} & প্রাণ দাও — জীবন্ত করো \\
\textbf{খালাকতুম (\arab{خلقتم})} & তোমরা সৃষ্টি করেছ \\
\hline
\end{tabular}
\end{table}

\#\#\#\# ৬. গুরুত্বপূর্ণ শিক্ষা (\lat{Key} \lat{Lessons})

\begin{enumerate}
\item \textbf{নিষেধের কারণ স্পষ্ট:} ছবি বানানোর দোষ এই — জীবনের সৃষ্টি আল্লাহর একান্ত বৈশিষ্ট্যের প্রতিদ্বন্দ্বিতা।
\item \textbf{পশু-পাখিও প্রাণবান:} পাখি-পশুর ছবি আঁকার সময়ও এই একই প্রশ্ন — প্রাণ দেওয়া সম্ভব? না। তাই সেগুলোও নিষিদ্ধ।
\item \textbf{মূর্তি আরও কঠোর:} দুই-মাত্রিক (আঁকা) ছবির চেয়ে মূর্তি (তিন-মাত্রিক) আরও কঠিন — কারণ তা আরও বেশি \lat{“}সৃষ্টির প্রতিদ্বন্দ্বী\lat{”}।
\end{enumerate}
\medskip\hrule\medskip

\subsubsection*{হাদিস ৩ — আয়েশার কুশন: ফেরেশতা ঢোকে না, ছবিওয়ালা ঘরে (সহীহ বুখারি ৫৯৫৭, ৫৯৬১; সহীহ মুসলিম ২১০৭)}

\#\#\#\# ১. সূত্র কার্ড

\begin{table}[h]\centering\small
\begin{tabular}{ll}
বিষয় & বিবরণ \\
\hline
\textbf{বর্ণনাকারী} & আয়েশা (রাঃ) — নবীজির স্ত্রী \\
\textbf{গ্রন্থ} & সহীহ বুখারি ৫৯৫৭, ৫৯৬১, ৫১৮১; সহীহ মুসলিম ২১০৭ \\
\textbf{অধ্যায়} & কিতাবুল লিবাস — বাবু মান কারিহাল কুউদা আলাস-সুওয়ারাহ (ছবির উপর বসতে অপছন্দ) \\
\textbf{সনদের মান} & \textbf{সহীহ} (বুখারি-মুসলিম) \\
\hline
\end{tabular}
\end{table}

\#\#\#\# ২. পটভূমি (\lat{Context})

আয়েশা (রাঃ) ছবিওয়ালা একটি \textbf{কুশন (\lat{numruqa})} কিনেছিলেন — যাতে নবীজি (সা.) বসে হেলান দিতে পারেন। নবীজি (সা.) তা দেখে দরজায় দাঁড়িয়ে থেকে ঘরে ঢুকলেন না। আয়েশা (রাঃ) বিচলিত হয়ে তাওবা করলেন এবং জানতে চাইলেন — কী অপরাধ করেছেন?

\#\#\#\# ৩. পূর্ণ বর্ণনা (বাংলা, \lat{pure})

\begin{quote}
আয়েশা (রাঃ) বলেন: আমি ছবিওয়ালা একটি কুশন কিনলাম। রাসূলুল্লাহ (সাল্লাল্লাহু আলাইহি ওয়া সাল্লাম) তা দেখে দরজায় দাঁড়িয়ে থেকে ঘরে প্রবেশ করলেন না। আমি তাঁর চেহারায় অপছন্দের লক্ষণ দেখে বললাম: হে আল্লাহর রাসূল! আমি আল্লাহ ও তাঁর রাসূলের কাছে তাওবা করছি — আমি কী অপরাধ করেছি? \\
তিনি বললেন: \textbf{\lat{“}এই কুশনের ব্যাপারটি কী?\lat{”}} আমি বললাম: এটা আপনার বসা ও হেলান দেওয়ার জন্য কিনেছি। \\
তিনি বললেন: \textbf{\lat{“}এই ছবিগুলোর মালিকদের কিয়ামতের দিন শাস্তি দেওয়া হবে, আর বলা হবে — যা সৃষ্টি করেছ তাতে প্রাণ দাও।\lat{”}} তিনি আরও বললেন: \textbf{\lat{“}নিশ্চয়ই যে ঘরে ছবি থাকে, সে ঘরে ফেরেশতারা প্রবেশ করেন না।\lat{”}}
\end{quote}

\#\#\#\# ৪. ধাপে ধাপে বিশ্লেষণ

\textbf{ধাপ ১ — নবীজির প্রতিক্রিয়া:} ছবিওয়ালা কুশন দেখে ঘরে ঢোকেননি — ছবির প্রতি তীব্র \textbf{\lat{displeasure} (অপছন্দ)} প্রকাশের প্রমাণ।

\textbf{ধাপ ২ — দুটি বিধান:} এই হাদিসে দুইটি বিধান — (ক) ছবি নির্মাতাদের শাস্তি, (খ) ছবি-ওয়ালা ঘরে ফেরেশতা প্রবেশ করে না।

\textbf{ধাপ ৩ — \lat{“}ছবির মালিক\lat{”} (\arab{أصحاب} \arab{هذه} \arab{الصور}):} শাস্তির বিবৃতি ছবি \textbf{নির্মাতাদের}; কিন্তু ফেরেশতা না ঢোকার বিবৃতি ঘরে ছবি \textbf{রাখার} ক্ষেত্রে। অর্থাৎ — আঁকা হারাম, আর ঘরে (অবজ্ঞা না করা) প্রাণীর ছবি রাখাও ফেরেশতাদের আগমন রোধ করে।

\#\#\#\# ৫. শব্দের ব্যাখ্যা (\lat{Vocabulary})

\begin{table}[h]\centering\small
\begin{tabular}{ll}
শব্দ & অর্থ \\
\hline
\textbf{নুমরুকাহ (\arab{نمرقة})} & কুশন/গদি — বসা ও হেলান দেওয়ার বালিশ \\
\textbf{তাসাবির (\arab{تصاوير})} & ছবি — এখানে প্রাণীর ছবি \\
\textbf{আসহাবু হাযিহিস-সুওয়ার} & এই ছবিগুলোর মালিক/নির্মাতা \\
\textbf{আল-মালায়িকাহ (\arab{الملائكة})} & ফেরেশতারা — রহমতের ফেরেশতা \\
\hline
\end{tabular}
\end{table}

\#\#\#\# ৬. গুরুত্বপূর্ণ শিক্ষা (\lat{Key} \lat{Lessons})

\begin{enumerate}
\item \textbf{ঘরে প্রাণীর ছবি রাখা:} ঝুলানো-সজ্জিত প্রাণীর ছবি (মানুষ/পশু-পাখি) ঘরে রাখা ফেরেশতার আগমন বাধাগ্রস্ত করে — সতর্কতা।
\item \textbf{কুশন-বালিশে ছবি:} এখানে ছবি আঁকার \textbf{নিষেধ} নিশ্চিত; তবে এ হাদিসের ধারাবাহিকতায় (মুসলিম ২১০৮) বালিশ-কুশনের মতো অবজ্ঞার বস্তুতে ছবির ক্ষেত্রে ব্যতিক্রমের আলোচনা আছে — পরের হাদিসে।
\item \textbf{তাওবার মানসিকতা:} আয়েশা (রাঃ) সঙ্গে সঙ্গে তাওবা করেছেন — ভুল বুঝলেই সংশোধন।
\end{enumerate}
\medskip\hrule\medskip

\subsubsection*{হাদিস ৪ — ফেরেশতা ঢোকে না: কুকুর বা ছবি-ওয়ালা ঘরে (সহীহ বুখারি ৩২২৬, সহীহ মুসলিম ২১০৬)}

\#\#\#\# ১. সূত্র কার্ড

\begin{table}[h]\centering\small
\begin{tabular}{ll}
বিষয় & বিবরণ \\
\hline
\textbf{বর্ণনাকারী} & আবু তালহা আনসারী (রাঃ); যায়িদ ইবনে খালিদ (রাঃ)-এর বর্ণনায় \\
\textbf{গ্রন্থ} & সহীহ বুখারি ৩২২৬, ৫৯৪৯; সহীহ মুসলিম ২১০৬; সুনানে আবু দাউদ ৪১৫৩; সুনানে তিরমিযী ২৮৪ \\
\textbf{অধ্যায়} & কিতাবু বাদয়িল খালক — ফেরেশতা ও ছবি/কুকুরের অধ্যায় \\
\textbf{সনদের মান} & \textbf{সহীহ} (বুখারি-মুসলিম) \\
\hline
\end{tabular}
\end{table}

\#\#\#\# ২. পটভূমি (\lat{Context})

এই হাদিসটি ছবি ও কুকুর — দুই জিনিসের কথা বলেছে, যার উপস্থিতি ফেরেশতাদের প্রবেশ বাধাগ্রস্ত করে। যায়িদ ইবনে খালিদ (রাঃ) আয়েশা (রাঃ)-এর কাছে এসে বিষয়টি যাচাই করেছিলেন — আয়েশা (রাঃ) নিজে শোনেননি, তবে তাঁর দেখা \textbf{পর্দার ঘটনা} (পরের হাদিস) বর্ণনা করেন।

\#\#\#\# ৩. পূর্ণ বর্ণনা (বাংলা, \lat{pure})

\begin{quote}
\textbf{\lat{“}ফেরেশতারা সে ঘরে প্রবেশ করেন না, যে ঘরে কুকুর বা ছবি (মূর্তি) থাকে।\lat{”}}
\end{quote}

\#\#\#\# ৪. ধাপে ধাপে বিশ্লেষণ

\textbf{ধাপ ১ — প্রাণীর ছবিই উদ্দেশ্য:} ইমাম খাত্তাবির ব্যাখ্যা — যে ছবির কারণে ঘরে ফেরেশতা ঢোকে না, তা হলো \textbf{প্রাণবান (জীবনবিশিষ্ট) প্রাণীর ছবি}, যার মাথা কাটা হয়নি এবং যা অবজ্ঞার বস্তু (বিছানা-বালিশ) নয়।

\textbf{ধাপ ২ — \lat{“}কুকুর\lat{”} ও \lat{“}ছবি\lat{”} একসাথে:} উভয়ই ফেরেশতার আগমনের বাধা — ছবির কারণ দাঁড় করানো হলো \lat{“}প্রাণীর ছবি\lat{”}।

\textbf{ধাপ ৩ — ব্যতিক্রমের ইঙ্গিত:} খাত্তাবির মতে — মাথাহীন, অসম্পূর্ণ বা অবজ্ঞার বস্তুতে থাকা ছবি ফেরেশতাদের বাধা নয়। এটা পরবর্তী হাদিসে (আয়েশার বালিশ) স্পষ্ট।

\#\#\#\# ৫. শব্দের ব্যাখ্যা (\lat{Vocabulary})

\begin{table}[h]\centering\small
\begin{tabular}{ll}
শব্দ & অর্থ \\
\hline
\textbf{কালব (\arab{كلب})} & কুকুর \\
\textbf{সুওরাহ (\arab{صورة})} & ছবি/মূর্তি — এখানে প্রাণীর ছবি \\
\textbf{মুমতাহানাহ (\arab{ممتهنة})} & অবজ্ঞার/ব্যবহৃত বস্তু — মাটিতে মোছা, বসার জিনিস \\
\hline
\end{tabular}
\end{table}

\#\#\#\# ৬. গুরুত্বপূর্ণ শিক্ষা (\lat{Key} \lat{Lessons})

\begin{enumerate}
\item \textbf{ঘরের পরিবেশ:} মুমিনের ঘরে প্রাণীর ঝুলানো ছবি ও অবজ্ঞা না করা ছবি না রাখার চেষ্টা — ফেরেশতাদের আগমনের পথ খোলা রাখা।
\item \textbf{দুটি শর্ত মিললে ছবি অ-বাধক:} মাথাহীন অথবা অবজ্ঞার বস্তুতে (বালিশ-কার্পেট) ছবি — পরবর্তী আলোচনায় বিস্তারিত।
\item \textbf{সতর্কতা ও ভারসাম্য:} নিষেধ মানতে গিয়ে বাড়াবাড়ি নয় — অবজ্ঞার বস্তুতে ছবির ব্যতিক্রম পরের হাদিসে।
\end{enumerate}
\medskip\hrule\medskip

\subsection*{পার্ট ২ — ব্যতিক্রম ও সীমারেখা}

\subsubsection*{হাদিস ৫ — ইবনে আব্বাসের উপদেশ: \lat{“}গাছের ছবি বানাও\lat{”} (সহীহ মুসলিম ২১১০খ, সহীহ বুখারি ৫৯৬৩)}

\#\#\#\# ১. সূত্র কার্ড

\begin{table}[h]\centering\small
\begin{tabular}{ll}
বিষয় & বিবরণ \\
\hline
\textbf{বর্ণনাকারী} & ইবনে আব্বাস (রাঃ) \\
\textbf{গ্রন্থ} & সহীহ মুসলিম ২১১০খ; সহীহ বুখারি ৫৯৬৩ \\
\textbf{অধ্যায়} & কিতাবুল লিবাস — ছবি নির্মাতাদের শাস্তি \\
\textbf{সনদের মান} & \textbf{সহীহ} (বুখারি-মুসলিম) \\
\hline
\end{tabular}
\end{table}

\#\#\#\# ২. পটভূমি (\lat{Context})

একজন পেশাদার ছবি-শিল্পী ইবনে আব্বাসের (রাঃ) কাছে এসে বললেন — তিনি এই ছবিগুলো আঁকেন এবং এটিই তার জীবিকা। ইবনে আব্বাস (রাঃ) নবীজির (সা.) বাণী শোনালেন। এই হাদিসটি সবচেয়ে স্পষ্ট প্রমাণ — \textbf{প্রাণবান সৃষ্টির ছবি নিষিদ্ধ, জড়বস্তুর ছবি জায়েজ}।

\#\#\#\# ৩. পূর্ণ বর্ণনা (বাংলা, \lat{pure})

\begin{quote}
এক ব্যক্তি ইবনে আব্বাসকে (রাঃ) বলল: আমি এমন এক ব্যক্তি, যে এই ছবিগুলো আঁকি। ইবনে আব্বাস (রাঃ) বললেন: আমি রাসূলুল্লাহকে (সাল্লাল্লাহু আলাইহি ওয়া সাল্লাম) বলতে শুনেছি: \\
\textbf{\lat{“}যে ব্যক্তি দুনিয়ায় ছবি বানাবে, কিয়ামতের দিন তাকে তাতে প্রাণ দেওয়ার ভার দেওয়া হবে, আর সে প্রাণ দিতে পারবে না।\lat{”}} \\
এরপর ইবনে আব্বাস (রাঃ) বললেন: \textbf{\lat{“}যদি তুমি ছবি আঁকতেই চাও, তাহলে গাছের ছবি আঁকো এবং এমন কোনো বস্তুর ছবি আঁকো যাতে প্রাণ নেই।\lat{”}}
\end{quote}

\#\#\#\# ৪. ধাপে ধাপে বিশ্লেষণ

\textbf{ধাপ ১ — প্রশ্নের স্পষ্টতা:} শিল্পী বললেন \lat{“}এই ছবিগুলো\lat{”} — অর্থাৎ প্রাণীর ছবি; ইবনে আব্বাস (রাঃ) নিষেধের হাদিস শোনালেন।

\textbf{ধাপ ২ — প্রাণ দেওয়ার ভার:} নিষেধের মূল কারণ আবারও — প্রাণ দেওয়া মানুষের সাধ্যের বাইরে; ছবি বানানো সেই অসম্ভব দাবির অনুকরণ।

\textbf{ধাপ ৩ — বৈধ বিকল্প:} \lat{“}গাছের ছবি বানাও\lat{”} — এই বাক্যই \textbf{জায়েজের সীমারেখা} নির্ধারণ করে: প্রাণহীন জড়বস্তু (গাছ, পাহাড়, নদী, বাড়ি, নৌকা, ফুল) আঁকা সম্পূর্ণ জায়েজ। এতেই প্রমাণ — নিষেধ \textbf{শুধু প্রাণবান সৃষ্টির}, মানুষ বা পশু-পাখি সবই এতে সমান।

\#\#\#\# ৫. শব্দের ব্যাখ্যা (\lat{Vocabulary})

\begin{table}[h]\centering\small
\begin{tabular}{ll}
শব্দ & অর্থ \\
\hline
\textbf{উসাওয়িরু (\arab{أصور})} & ছবি আঁকি \\
\textbf{আন-নাফখ ফীহির রূহ} & তাতে প্রাণ (আত্মা) ফুঁকে দেওয়া \\
\textbf{শাজারাহ (\arab{شجرة})} & গাছ \\
\textbf{মা- লা রূহা ফিহ (\arab{ما} \arab{لا} \arab{روح} \arab{فيه})} & যাতে প্রাণ নেই — জড়বস্তু \\
\hline
\end{tabular}
\end{table}

\#\#\#\# ৬. গুরুত্বপূর্ণ শিক্ষা (\lat{Key} \lat{Lessons})

\begin{enumerate}
\item \textbf{প্রাণবান = নিষিদ্ধ, জড় = জায়েজ:} এই এক বাক্যে দ্বিধার সমাধান — মানুষ, পশু, পাখি (সব প্রাণবান) হারাম; গাছ-পাহাড়-বাড়ি-নৌকা (জড়বস্তু) জায়েজ।
\item \textbf{পেশার জন্য জায়েজ নয়:} জীবিকার জন্য ছবি আঁকলেও নিষেধ বহাল — কারণ হাদিসে শিল্পী নিজের জীবিকার কথা বলেই শাস্তির কথা শুনেছেন।
\item \textbf{বিকল্প আছে:} ইসলাম শিল্প নিষিদ্ধ করে না — শুধু প্রাণবান সৃষ্টির অনুকরণ থেকে বিরত রাখে; প্রকৃতি-জড়বস্তুর শিল্পে খোলা দরজা।
\end{enumerate}
\medskip\hrule\medskip

\subsubsection*{হাদিস ৬ — আয়েশার পর্দা: ছবি ছিঁড়ে বালিশ বানানো (সহীহ মুসলিম ২১০৮, সহীহ বুখারি ৫৯৫৪)}

\#\#\#\# ১. সূত্র কার্ড

\begin{table}[h]\centering\small
\begin{tabular}{ll}
বিষয় & বিবরণ \\
\hline
\textbf{বর্ণনাকারী} & আয়েশা (রাঃ) \\
\textbf{গ্রন্থ} & সহীহ মুসলিম ২১০৮; সহীহ বুখারি ৫৯৫৪; সহীহ বুখারি ২৪৮৭ \\
\textbf{অধ্যায়} & কিতাবুল লিবাস — অবজ্ঞার বস্তুতে ছবি \\
\textbf{সনদের মান} & \textbf{সহীহ} (বুখারি-মুসলিম) \\
\hline
\end{tabular}
\end{table}

\#\#\#\# ২. পটভূমি (\lat{Context})

আয়েশা (রাঃ) দরজায় ছবিওয়ালা একটি \textbf{পর্দা (\lat{qiram})} লাগিয়েছিলেন। নবীজি (সা.) সফর থেকে ফিরে তা দেখে অপছন্দ প্রকাশ করেন এবং বলেন — আল্লাহ আমাদেরকে পাথর-মাটিকে (এভাবে) পোশাক পরাতে বলেননি। আয়েশা (রাঃ) পর্দাটি কেটে দুইটি বালিশ (কুশন) বানালেন — তাতে নবীজি (সা.) কোনো আপত্তি করেননি।

\#\#\#\# ৩. পূর্ণ বর্ণনা (বাংলা, \lat{pure})

\begin{quote}
আয়েশা (রাঃ) বলেন: আমি দরজায় একটি পর্দা লাগালাম, যাতে ছবি ছিল। রাসূলুল্লাহ (সাল্লাল্লাহু আলাইহি ওয়া সাল্লাম) সফর থেকে ফিরে তা দেখে চেহারায় অসন্তুষ্টির আলামত ফুটে উঠল। তিনি তা টেনে নামালেন, এমনকি ছিঁড়ে ফেললেন বা টুকরা করে ফেললেন, আর বললেন: \\
\textbf{\lat{“}আল্লাহ আমাদেরকে পাথর বা মাটিকে পোশাক পরানোর নির্দেশ দেননি।\lat{”}} \\
আয়েশা (রাঃ) বলেন: আমরা পর্দাটি কেটে দুইটি বালিশ বানালাম, ভেতরে খেজুরের আঁশ ভরে দিলাম — তাতে তিনি আমাকে দোষারোপ করেননি। \\
অন্য বর্ণনায়: তিনি বললেন — \textbf{\lat{“}আয়েশা! কিয়ামতের দিন সবচেয়ে কঠিন শাস্তি হবে তাদের, যারা আল্লাহর সৃষ্টির অনুকরণ করে।\lat{”}}
\end{quote}

\#\#\#\# ৪. ধাপে ধাপে বিশ্লেষণ

\textbf{ধাপ ১ — সম্মানিত স্থানে ছবি:} দেয়ালে-দরজায় ঝোলানো, উপরে-মাথার দিকে থাকা ছবি — সর্বোচ্চ নিষেধ; নবীজি (সা.) এত তীব্র প্রতিক্রিয়া দেখালেন।

\textbf{ধাপ ২ — পর্দা ছিঁড়ে ফেলা:} ঝুলানো ছবি \textbf{ধ্বংস/সংশোধন} করা দরকার — এখানে নবীজি (সা.) নিজেই টেনে ছিঁড়লেন।

\textbf{ধাপ ৩ — বালিশে ছবি অনুমোদিত:} পর্দাটি বালিশে রূপান্তর — এখন সেটি \textbf{মোছা-ব্যবহৃত (\lat{mumtahana})} বস্তু; বসা-হেলান দেওয়ার জিনিস; এ অবস্থায় নবীজি (সা.) আপত্তি করেননি। এটিই ব্যতিক্রমের দলিল — \textbf{অবজ্ঞার/ব্যবহৃত জিনিসে (বালিশ, কুশন, কার্পেট, মুছার কাপড়) প্রাণীর ছবি জায়েজ} (অধিকাংশ আলেমের মতে)।

\#\#\#\# ৫. শব্দের ব্যাখ্যা (\lat{Vocabulary})

\begin{table}[h]\centering\small
\begin{tabular}{ll}
শব্দ & অর্থ \\
\hline
\textbf{কিরাম (\arab{قرام})} & পর্দা — দরজায় ঝোলানো ছবিওয়ালা চাদর \\
\textbf{তুসাব্বিরুন (\arab{تصاوير})} & ছবি — প্রাণীর ছবি \\
\textbf{মুমতাহানাহ (\arab{ممتهنة})} & অবজ্ঞার/ব্যবহৃত বস্তু — যার প্রতি সম্মান নেই \\
\textbf{উসাদাতান (\arab{وسادتان})} & দুইটি বালিশ \\
\hline
\end{tabular}
\end{table}

\#\#\#\# ৬. গুরুত্বপূর্ণ শিক্ষা (\lat{Key} \lat{Lessons})

\begin{enumerate}
\item \textbf{ঝুলানো/টাঙানো প্রাণীর ছবি — সবচেয়ে কঠিন অবস্থা:} দেয়ালে-দরজায়, সম্মানের স্থানে প্রাণীর ছবি রাখা ফেরেশতার আগমন রোধ ও নবীজির (সা.) তীব্র অপছন্দের বিষয়।
\item \textbf{ব্যবহৃত জিনিসে ছবি — ব্যতিক্রম:} বালিশ, কুশন, কার্পেট, গদি — যেখানে বসা-মোছা হয় — সেখানে ছবি অধিকাংশ আলেমের মতে জায়েজ; এখানেই দ্বিধা-সমাধানের ভারসাম্য।
\item \textbf{সংশোধনের পথ:} ঝুলানো ছবি টাঙানোর বদলে ব্যবহার্য বস্তুতে রূপান্তর — নবীজির (সা.) দেখানো সমাধান।
\end{enumerate}
\medskip\hrule\medskip

\subsection*{পার্ট ৩ — হাদিস থেকে মূল উত্তর (\lat{The} \lat{Core} \lat{Answer})}

\subsubsection*{প্রশ্ন: শুধু মানুষের ছবি হারাম, নাকি পশুপাখিরও?}

\textbf{উত্তর: পশুপাখির ছবিও হারাম।} হাদিসগুলোর ভাষা থেকে এটি পরিষ্কার:

\begin{itemize}
\item হাদিস ১-২: \lat{“}ছবি নির্মাতারা\lat{”}, \lat{“}যারা ছবি বানায়\lat{”} — শব্দ \textbf{সাধারণ}, মানুষে সীমাবদ্ধ নয়।
\item হাদিস ৫: ইবনে আব্বাস (রাঃ) স্পষ্ট বলেন — \lat{“}প্রাণহীন বস্তুর ছবি বানাও\lat{”}; অর্থাৎ \textbf{যেকোনো প্রাণবান সৃষ্টির} ছবি নিষিদ্ধ — মানুষ, পশু, পাখি, মাছ, কীট — সব।
\item ইমাম নববির (রহ.) বক্তব্য — ছবি বলতে সব \textbf{প্রাণবান সৃষ্টির ছবি} উদ্দেশ্য; প্রাণহীন (গাছ, পাহাড়, নদী) জায়েজ।
\item আলেমদের ঐকমত্য — এই বিষয়ে মূলত কোনো মতভেদ নেই: \textbf{প্রাণবান সৃষ্টির ছবি আঁকা/গড়া নিষিদ্ধ; জড়বস্তুর ছবি জায়েজ।} মতভেদ শুধু ফটোগ্রাফি, ডিজিটাল ও \lat{AI}-এর মতো আধুনিক মাধ্যম নিয়ে — সেগুলো ব্যবহারিক গাইড ফাইলে।
\end{itemize}
\subsubsection*{মাযহাবগুলোর অবস্থান (সংক্ষিপ্ত)}

\begin{table}[h]\centering\small
\begin{tabular}{llll}
মাযহাব & প্রাণবান সৃষ্টির ছবি (আঁকা/গড়া) & জড়বস্তুর ছবি & অবজ্ঞার বস্তুতে (বালিশ-কার্পেট) ছবি \\
\hline
\textbf{হানাফি} & নিষিদ্ধ (সম্পূর্ণ, মাথাওয়ালা) & জায়েজ & জায়েজ \\
\textbf{শাফেয়ি} & নিষিদ্ধ (প্রসিদ্ধ মত) & জায়েজ & জায়েজ \\
\textbf{মালেকি} & নিষিদ্ধ/মাকরূহ (মতভেদ) & জায়েজ & জায়েজ (ব্যবহৃত) \\
\textbf{হাম্বলি} & নিষিদ্ধ & জায়েজ & জায়েজ (মোছার জিনিস) \\
\hline
\end{tabular}
\end{table}

\begin{quote}
দ্রষ্টব্য: খাত্তাবি (রহ.)-এর মতো কিছু শাফেয়ি আলেম দুই-মাত্রিক ছবির ক্ষেত্রে হালকা মত দিয়েছেন; তবে \textbf{জমহুর (সংখ্যাগরিষ্ঠ) আলেমের প্রসিদ্ধ ও নিরাপদ মত — প্রাণবান সৃষ্টির ছবি তৈরি নিষিদ্ধ}। হানাফিদের মতে মাথাবিহীন/অসম্পূর্ণ ছবি জায়েজ — এটি একটি পরিচিত ইখতিলাফ, এখানে উল্লেখ করা হলো সততার সাথে।
\end{quote}

\medskip\hrule\medskip

\subsection*{রেফারেন্স তালিকা}

\begin{itemize}
\item সহীহ বুখারি: ৫৯৫০, ৫৯৫১, ৫৯৫৪, ৫৯৫৭, ৫৯৬১, ৫৯৬৩, ৫৯৪৯, ৩২২৬, ৫১৮১, ২৪৮৭
\item সহীহ মুসলিম: ২১০৬, ২১০৭, ২১০৮, ২১০৯, ২১১০খ
\item সুনানে আবু দাউদ: ৪১৫৩; সুনানে তিরমিযী: ২৮৪
\item মিশকাতুল মাসাবীহ: ৪৪৮৯ (তাহকীকসহ বর্ণনা)
\item ব্যাখ্যা: ফাতহুল বারি, মিরকাতুল মাফাতীহ, শারহু নববী (সাহিহ মুসলিমের ব্যাখ্যা), আল-আজকিযাহ (খাত্তাবি)
\end{itemize}

\end{document}
